% !TEX root = ../algebra_02.tex

\section{Строение простых алгебраических расширений. Присоединение к полю корня неприводимого многочлена}

\begin{Theorem}{Об изоморфизме простого расширения}{}
	Пусть $L = K(x)$ --- простое расширение, где элемент $x \in L$ алгебраичен над $K$. Пусть $f_0 = \operatorname{Irr}_K(x)$ --- его минимальный многочлен, а $\deg f_0 = d$.
	Тогда существует изоморфизм полей:
	\[ K[t]/(f_0) \xrightarrow{\sim} L \]
	При этом элементы смежных классов вида $\alpha_0 + \alpha_1 t + \dots + \alpha_{d-1} t^{d-1}$ переходят в элементы $\alpha_0 + \alpha_1 x + \dots + \alpha_{d-1} x^{d-1} \in L$.
\end{Theorem}

\begin{proof}
	Рассмотрим отображение вычисления многочлена в точке $x$:
	\[ \varphi: K[t] \to L, \quad \sum_{i=0}^n \alpha_i t^i \mapsto \sum_{i=0}^n \alpha_i x^i \]
	Очевидно, что $\varphi$ является гомоморфизмом колец.
	Найдем ядро этого гомоморфизма. По определению, многочлен $f(t)$ зануляется в точке $x$ тогда и только тогда, когда он делится на минимальный многочлен $f_0$. Следовательно, $\ker \varphi = (f_0)$.

	По теореме о гомоморфизме колец, индуцированное отображение $\tilde{\varphi}$ задаёт изоморфизм:
	\[ K[t]/(f_0) \xrightarrow{\sim} \operatorname{Im} \varphi \]
	Поскольку минимальный многочлен $f_0$ неприводим, идеал $(f_0)$ максимален, а значит, факторкольцо $K[t]/(f_0)$ является полем. Следовательно, образ $\operatorname{Im} \varphi$ также является полем.

	Заметим, что:
	\begin{itemize}
		\item $\operatorname{Im} \varphi \supset K$, так как для скаляров $\tilde{\varphi}(\alpha) = \alpha$;
		\item $\operatorname{Im} \varphi \ni x$, так как $\tilde{\varphi}(\bar{t}) = x$;
		\item $\operatorname{Im} \varphi \subset L$ по построению.
	\end{itemize}
	Так как $L = K(x)$ --- наименьшее поле, содержащее $K$ и $x$, а $\operatorname{Im} \varphi$ --- поле с такими же свойствами, получаем $\operatorname{Im} \varphi = L$.
\end{proof}

\begin{Definition}{$K$-изоморфизм расширений}{}
	Пусть даны расширения $L/K$ и $L'/K$. Изоморфизм полей $\varphi: L \to L'$ называется \textbf{$K$-изоморфизмом расширений}, если он оставляет все элементы базового поля $K$ на месте:
	\[ \forall a \in K: \varphi(a) = a \]
\end{Definition}

Таким образом, доказанная теорема показывает, что простое алгебраическое расширение $K(x)$ единственно с точностью до $K$-изоморфизма и полностью определяется минимальным многочленом элемента $x$.

\begin{Theorem}{О существовании простого расширения}{}
	Пусть $f_0 \in K[t]$ --- неприводимый многочлен.
	Тогда существует расширение $L/K$ и элемент $x \in L$ такие, что $L = K(x)$ и $\operatorname{Irr}_K(x) = f_0$.
\end{Theorem}

\begin{proof}
	Сконструируем поле $L = K[t]/(f_0)$. Мы можем рассматривать $L$ как расширение $K$, отождествляя каждый элемент $\alpha \in K$ с классом смежности скалярного многочлена $\bar{\alpha} \in K[t]/(f_0)$.
	Обозначим класс смежности переменной $t$ за $x = \bar{t}$. Тогда любой элемент из $L$ выражается как многочлен от $x$, то есть $L = K(x)$.

	Покажем, что $x$ --- корень $f_0$. Если $f_0(t) = \sum \alpha_i t^i$, то:
	\[ f_0(x) = \sum \alpha_i \bar{t}^i = \overline{\sum \alpha_i t^i} = \overline{f_0(t)} = 0 \]
	Таким образом, $f_0 \in \operatorname{Ann}(x)$.
	При этом для любого многочлена $f \in K[t]$ выполняется:
	\[ f(x) = 0 \iff \overline{f(t)} = 0 \iff f(t) \in (f_0) \iff f_0 \mid f \]
	Следовательно, $f_0$ --- многочлен наименьшей степени, зануляющий $x$, то есть $\operatorname{Irr}_K(x) = f_0$.
\end{proof}

\begin{Corollary}{}{}
	Если элемент $x$ алгебраичен над $K$, то степень порожденного им простого расширения равна степени его минимального многочлена:
	\[ [K(x):K] = \deg \operatorname{Irr}_K(x) \]
\end{Corollary}

\begin{Corollary}{}{}
	Пусть $L/K$ --- конечное расширение, а $x \in L$.
	Тогда степень алгебраического элемента $x$ делит степень расширения:
	\[ \deg_K x \mid [L:K] \]
\end{Corollary}

\begin{proof}
	Рассмотрим башню полей $K \subset K(x) \subset L$. По теореме о мультипликативности степени конечных расширений:
	\[ [L:K] = [L:K(x)] \cdot [K(x):K] \]
	Подставляя $[K(x):K] = \deg_K x$, получаем требуемое утверждение.
\end{proof}
